\section{Actual up and down neighbors and their chart naturality}
\label{sec:neighbor-geometry}

\status{NB1--NB5 construct and identify the actual neighboring decorated
lattices, their residue maps and their chart labels. NB6--NB8 prove the
full up/down naturality statements using AD3's supported splice. The
initial adelic realization of a valid decorated point remains an explicit
input. No equality of integral lattices is deduced merely from equal rank.}

\paragraph{Conventions and exact inputs.}
Put $\mathcal O=\mathcal O_E$, $V=E^n$, and
$h(x,y)=\bar x^{\mathsf T}y$. For a full $\mathcal O$-lattice $M$ write
\[
 M^\#=\{y\in V:h(M,y)\subseteq\mathcal O\},\qquad
 M_P=\mathcal O_{(P)}M\subset V.
\]
Here $M_P$ is localization inside the rational vector space; completed
lattices appear only when applying AD3. All sums and ideal products
below are sums and products of the actual integral submodules of $V$.
Let $d=(N,(f_Q)_{Q\in S_T},(\mathrm{be}_R)_{R\in S_G})$ be a valid
source decoration. Thus $N=N^\#$ is full, each $f_Q$ is a nonzero
residue-linear map $N_Q\to k_Q=\mathcal O/Q$, extended by zero outside
$N_Q$, and each $\mathrm{be}_R$ gives the prescribed residue basis
isomorphism at its auxiliary prime. Put
\[
 L=\{x\in N:f_Q(x)=0\text{ for all }Q\in S_T\},\qquad H=L^\#.
\]
Fix $P\in S_T$ such that $P\ne\bar P$, no member of $S_T$ is $\bar P$,
and all auxiliary and twist-frame primes avoid $P,\bar P$. These are
the precise \texttt{MoveHyp} and frame-disjointness requirements.
For an up move take $q\in P\setminus P^2$. In the true-column application
$q=T_u$: FC1 gives $v_P(\pi_u)=1$, $v_P(\bar\pi_u)=0$, and
$T_u=\pi_u\bar\pi_u$, so this is the actual rational integer $T_u$,
not a replacement uniformizer.

For the actual data the separation at the moved prime follows from
the stated prime and twist laws. No selected type prime can be the
conjugate of another, because $(\pi_u,\bar\pi_u)=\mathcal O$.
The auxiliary rational primes are different from $T_u$. A twist
prime has residue characteristic outside the set $S$ containing
$T_u$, so it cannot contain $T_u$; restriction to $\mathcal O_F$
proves the same assertion for every prime above the moved base-field
prime. Thus all auxiliary and twist primes avoid $P,\bar P$.
When using AD3 we retain its full structural hypotheses on these
actual sets: pairwise disjointness and absence of a conjugate pair
in their union. These are supplied by the same prime/selection/twist
data used in the chart construction, rather than by any neighbor
existence assertion.

The application in NB5--NB8 retains the supplied
\texttt{TC.Ok}, \texttt{TC.Pins}, $l:\texttt{I.Laws}$ and
$Z:\texttt{TC.Later}$, and the exact chart data
\texttt{ReadsK}, \texttt{IdY}, \texttt{IdV}, \texttt{IdL}.
FC constructs the indicated chart and projection data under its stated
geometric inputs. NB uses neither the final Forms0 assembly nor its
up/down clauses as a premise.

\paragraph{NB1: localization of the Hermitian dual uses the conjugate prime.}
For every full lattice $M$ one has
\[
 (M^\#)_P=\{y\in V:h(x,y)\in\mathcal O_{(P)}
                         \text{ for all }x\in M_{\bar P}\}.
 \tag{NB.1}
\]
Indeed, if $sy\in M^\#$ with $s\notin P$, and $tx\in M$ with
$t\notin\bar P$, then
$h(x,y)=h(tx,sy)/(\bar t s)$ is integral at $P$.
Conversely, take finitely many $\mathcal O$-generators $m_j$ of $M$.
The proposed condition puts every $h(m_j,y)$ in $\mathcal O_{(P)}$.
A product $s\notin P$ of their denominators gives
$h(M,sy)\subseteq\mathcal O$, and hence $sy\in M^\#$.

The pairing between $\mathcal O_{(\bar P)}^n$ and
$\mathcal O_{(P)}^n$ is perfect. In local bases its dual is given by
the inverse conjugate transpose of the basis matrix; taking this
operation twice returns the original local lattice. Thus
$M^{\#\#}=M$, first locally by this matrix calculation and then globally
by LP1. The same calculation proves that $M^\#$ is a full lattice:
clearing the finitely many basis-matrix denominators places it between
two nonzero scalar multiples of $\mathcal O^n$.
Consequently equality of the localized lattices at every finite prime
proves equality of the global integral lattices. Conjugation has been
used as an isomorphism between the two local rings; it is not being
treated as an automorphism of $\mathcal O_{(P)}$ when $P\ne\bar P$.
Equality $M_P=M'_P$ also gives the source's \texttt{SameAt} witnesses:
clear denominators outside $P$ on finitely many generators and multiply
them to obtain one $a\notin P$ with $aM\subseteq M'$, and repeat in
the opposite direction.

\paragraph{NB2: the actual up lattice and its residue kernel.}
Choose $x\in N$ with $f_P(x)=1$ and define the source's actual lattices
\[
 A=N+P^{-1}x,\qquad U=A^\#+\bar P A,
 \qquad P^{-1}=\{a\in E:aP\subseteq\mathcal O\}.
 \tag{NB.2}
\]
Here $P^{-1}x$ denotes its $\mathcal O$-span. Since
$N\subseteq A\subseteq P^{-1}N$, $A$ is full and
$\bar P N\subseteq A^\#\subseteq N$. Localizing the invertible ideal
$P^{-1}$, and using (NB.1), gives the complete table
\[
\begin{array}{c|ccc}
 &P&\bar P&R\notin\{P,\bar P\}\\ \hline
 A&N_P+q^{-1}\mathcal O_{(P)}x&N_{\bar P}&N_R\\
 A^\#&N_P&(A^\#)_{\bar P}&N_R\\
 U&A_P&(A^\#)_{\bar P}&N_R.
\end{array}                                                    \tag{NB.3}
\]
For example $\bar P$ is a unit ideal at $P$, while
$\bar P N_{\bar P}\subseteq(A^\#)_{\bar P}$; these facts justify
both exceptional entries in the last row. At $P$ and $\bar P$ these
entries are dual to each other by (NB.1); at every other prime they are
the entries of $N=N^\#$. Therefore $U=U^\#$ globally. The definition
also gives $U\subseteq A$ and $\bar P A\subseteq U$. This proves all
the lattice assertions of \texttt{upN\_spec}, including actual equality
away from the conjugate pair.

Update only the selected type functional, setting
\[
 f_P^+(y)=f_P(qy)\quad(y\in U_P),\qquad
 f_P^+(y)=0\quad(y\notin U_P).                       \tag{NB.4}
\]
Every $y\in U_P$ has the form $n+q^{-1}a x$ with
$n\in N_P$ and $a\in\mathcal O_{(P)}$. Residue-linearity extends to
this local ring by clearing a denominator outside $P$, so
$f_P^+(y)=a\bmod P$. Hence
\[
 \ker(f_P^+|_{U_P})=N_P,\qquad U_P/N_P\simeq k_P.     \tag{NB.5}
\]
In particular the new functional is nonzero on the global lattice:
multiply $q^{-1}x\in U_P$ by a denominator $s\notin P$ placing it in
$U$; its new value is the nonzero residue of $s$.

Keep every other $f_Q$ and every $\mathrm{be}_R$ literally unchanged.
Their local domains agree by (NB.3), so their local laws persist.
The global nonzero and surjectivity witnesses are obtained by clearing
a denominator outside their own prime and rescaling by its inverse
residue. The auxiliary kernel assertion persists as well: if
$z\in U$ has zero auxiliary residue at $R$, then $z\in R U_R$;
choose $s\notin R$ with $sz\in R U$ and $t$ with
$ts\equiv1\pmod R$. The equality
$z=t(sz)+(1-ts)z$ puts $z$ in $RU$. Thus this is a valid decoration,
the literal \texttt{d.up P q x}, with no fallback case.

For completeness the finite residue conditions give
\[
 L_P=\ker(f_P|_{N_P}),\qquad L_R=N_R\ (R\notin S_T),
 \qquad H_P=N_P.                                     \tag{NB.6}
\]
To obtain the first equality, clear a denominator into $N$, then
multiply by elements of each $Q\in S_T\setminus\{P\}$ which lie in
$Q$ but not in $P$; the other conditions vanish without changing the
localization at $P$. The second equality has the same proof. For the
third, $\bar P\notin S_T$ and (NB.1) apply. Writing $L^+,H^+$ for the
new decoration, (NB.3)--(NB.6) yield
\[
 N\subseteq H,\qquad (L^+)_P=N_P=H_P,\qquad
 (H^+)_P=U_P=(H+\mathcal Oq^{-1}x)_P.                \tag{NB.7}
\]
Together with (NB.4), these are the complete
\texttt{MoveUpLat.lattices\_at\_up} assertions.

\paragraph{NB3: the actual down lattice and its extended functional.}
Put
\[
 K_P^{\rm gl}=\{x\in N:f_P(x)=0\},\qquad
 D=K_P^{\rm gl}+P(K_P^{\rm gl})^\# .                 \tag{NB.8}
\]
The notation $K_P^{\rm gl}$ denotes a global lattice, not a localization.
It lies between $PN$ and $N$, so it is full. At every $R\ne P$ it has
localization $N_R$: multiply an element of $N$ by an element of
$P\setminus R$ to kill its residue. Its localization at $P$ is
$\ker(f_P|_{N_P})$. The paired-dual calculation in NB1 now gives
\[
 D_P=(K_P^{\rm gl})_P=L_P,\qquad
 D_{\bar P}=((K_P^{\rm gl})^\#)_{\bar P},\qquad
 D_R=N_R\quad(R\notin\{P,\bar P\}).                 \tag{NB.9}
\]
Indeed $((K_P^{\rm gl})^\#)_P=N_P$ and
$PN_P\subseteq(K_P^{\rm gl})_P$; at $\bar P$, the ideal $P$ is a unit.
The two exceptional entries of (NB.9) are dual to one another.
Thus $D=D^\#$, which proves the actual \texttt{dnN\_spec}.
Also $N_P/D_P\simeq k_P$, by the nonzero functional $f_P$.
These are local indices; $D$ and $N$ need not be comparable globally.

Let $g:D/PD\to k_P$ be any nonzero $k_P$-linear functional. For
$y\in D_P$, choose $s\notin P$ with $sy\in D$ and define
\[
 f_P^-(y)=\rho_P(s)^{-1}g([sy]),                     \tag{NB.10}
\]
where $\rho_P:\mathcal O\to k_P$ is the actual quotient map;
put $f_P^-(y)=0$ outside $D_P$. If $t$ is another denominator,
linearity gives $\rho_P(t)g([sy])=\rho_P(s)g([ty])$.
This proves independence, additivity and residue-linearity. On $D$
the formula is exactly $g\circ[\cdot]$, so it is nonzero there.
This is precisely both \texttt{TDat.extF} and \texttt{modP.ext};
no different quotient functional has been chosen.

Update $f_P$ to (NB.10), leave all other functionals and auxiliary maps
unchanged, and use the same witness and kernel arguments as in NB2.
The result is the actual valid \texttt{d.dn P (TDat.extF D P g)}.
If $L^-,H^-$ are its kernel lattice and Hermitian dual, then
\[
                         (H^-)_P=D_P=L_P.           \tag{NB.11}
\]
This follows from (NB.6), applied to the new valid decoration, and is
the substantive local identity in \texttt{dn\_locP\_lat}.

\paragraph{NB4: the actual moves preserve frames, coupling and chart label.}
The source's framed point move \texttt{PtY.mv} returns its input if the
proposed new lattice invalidates a frame or its coupling set. That
alternative does not occur here. Every twist frame is a basis of the
same localized lattice before and after (NB.3) or (NB.9). The actual
coupling sets depend only on the auxiliary and twist localizations:
at the penultimate step they are cosets defined from those local
lattices and the prescribed uniformizer powers, and at other steps
the coupling set is empty. FC2 described the reference case with no
twist primes. These localizations are all
unchanged. Thus the actual moved point has exactly the new decoration
of NB2 or NB3, the old rational frames, and the old coupling subset
$Y\subset V$. This proves the full content of \texttt{mv\_up} and
\texttt{mv\_dn}, including validity.

Here is why the same statement survives passage to the preferred
class representatives. Two framed points will be called equal on
the coupling data when they have the same twist frames, the same
$Y$, and the same auxiliary local lattices and maps
$\mathrm{be}_R$. Their chosen auxiliary lift-frames are then equal,
since those choices depend on the same lattice and reduction map.
Consequently their \texttt{cosetIn} labels are equal. Applying the
same column-group element to both points preserves this relation:
localization commutes with its $E$-linear rational part, the two
coupling sets have the same image, and the finite components transform
both frame and residue maps by the same formulas. This proves the
source's \texttt{SameG.smul} assertion directly.

At the penultimate step the preferred representative is chosen to
have the selected value of \texttt{cosetIn} from its nonempty set of
attainable values along the orbit. The old point and its raw neighbor
have the same value after every common group action; hence their two
sets of attainable values are identical. The fixed choice on this
set gives the same preferred value. At every other step both labels
are empty. The column index $b$ is unchanged, and \texttt{ReadsK.key}
identifies a chart label exactly by $b$ and this preferred coupling
value. Therefore both actual up and down moves keep \texttt{chOf}.
This proves \texttt{up\_keeps\_chart} and \texttt{dn\_keeps\_chart};
equality of the raw labels alone would not have established this.

We shall also use invariance under a pure rational unitary action.
At a twist prime the rational frame changes from $F_R$ to $\gamma F_R$,
so its coordinates of $\gamma y$ are identical to the old coordinates
of $y$. At an auxiliary prime, let $F_R,F'_R$ be the chosen old and
new lift-frames. Then
\[
 B_R=(F'_R)^{-1}\gamma F_R\in\mathrm{GL}_n(\mathcal O_{(R)}),
 \qquad B_R\equiv I\pmod R.                          \tag{NB.12}
\]
The reduction identity follows because both frames lift the same
transported auxiliary basis. The actual coupling law gives
$q_R F_R^{-1}y\in\mathcal O_{(R)}^n$, where the auxiliary rational
prime $q_R$ is a uniformizer. Thus
$(B_R-I)F_R^{-1}y\in\mathcal O_{(R)}^n$. The new coordinates and old
coordinates define exactly the same coupling coset. At twist primes
they were equal already, including any prescribed higher congruence
power. This proves \texttt{cosetIn\_smul\_unitary} with the actual
chosen lift-frames and supplies the rational-action condition needed
for the level-reading argument below.

\paragraph{NB5: transport by the actual moving element and compare projections.}
Let $c'$ be the class of either raw neighbor and let
$g_0=(\gamma,h_0)$ be the source's chosen moving element, so
\[
                      g_0\,r(c')=p_{\rm raw}.        \tag{NB.13}
\]
Its rational part $\gamma$ is unitary. The finite part changes only
the finite decoration; hence $\gamma$ carries the preferred new
$N,L,H$ to the corresponding raw $N,L,H$. For $L$ this follows by
pulling back all type functionals, and for $H$ it follows from
$h(\gamma x,\gamma y)=h(x,y)$. The rational matrix is therefore
exactly \texttt{gam(mv)}, not an independently chosen lattice isometry.

For an up move start with $z\in H_c$ satisfying
$\varphi_c(\operatorname{pr}_c z)=1$, where $\varphi_c$ is the
functional on $N_c/PN_c$. Choose the source's quotient representative
$x_0\in N_c$ of $\operatorname{pr}_c z$. FC5 gives
$z-x_0\in PH_c$ and $f_P(x_0)=1$. Since $q$ is a uniformizer,
$q^{-1}(z-x_0)\in(H_c)_P$. Thus (NB.7) and (NB.13) give
\[
 (\gamma L_{c'})_P=(H_c)_P,\qquad
 (\gamma H_{c'})_P=(H_c+\mathcal Oq^{-1}z)_P.         \tag{NB.14}
\]
For $x\in H_{c'}$ with $q\gamma x\in H_c$, choose a common
$s\notin P$ with $sx\in N_{c'}$ and $sq\gamma x\in N_c$.
Such a denominator exists by $(H_c)_P=(N_c)_P$ and its new analogue.
On these global elements (NB.4) and (NB.13) give
$f_{c',P}(sx)=f_{c,P}(sq\gamma x)$. FC5 identifies the two global
reductions with the projections. Their residue-linearity and
$\rho_P(s)\ne0$ then give the exact identity
\[
 \varphi_{c'}(\operatorname{pr}_{c'}x)
       =\varphi_c(\operatorname{pr}_c(q\gamma x)).    \tag{NB.15}
\]

For a down move, let $g:D_c/PD_c\to k_P$ be its specified nonzero
functional and put $K_c=L_c+PH_c$. Its localization is $(D_c)_P$
by (NB.6) and (NB.9). Equations (NB.11) and (NB.13) give
\[
                         (\gamma H_{c'})_P=(L_c)_P. \tag{NB.16}
\]
If $x\in H_{c'}$ and $\gamma x\in K_c$, choose $s\notin P$ with
$sx\in N_{c'}$. Equation (NB.13) places $s\gamma x$ in the actual
global lattice $D_c$. By (NB.10) its new residue is $g([s\gamma x])$.
FC5 identifies $[sx]$ and $[s\gamma x]$ with
$\operatorname{pr}_{c'}(sx)$ and $\operatorname{prk}_c(s\gamma x)$.
Cancel the same nonzero residue $\rho_P(s)$ to obtain
\[
 \varphi_{c'}(\operatorname{pr}_{c'}x)
                   =g(\operatorname{prk}_c(\gamma x)).           \tag{NB.17}
\]
This proves \texttt{dn\_level\_functional} for the exact source
projection maps and chosen moving matrix. No assertion that $K_c=D_c$
globally has been used; only their equality at $P$ is required.

\paragraph{NB6: a supported adelic factorization with the exact character.}
Assume the retained realization input: each valid framed decorated
point under consideration has a unitary adelic realizer in the full
\texttt{IsPtOf} sense. Take $h=y_ck$ realizing $r(c)$ and a realizer
of the actual raw neighbor. NB2--NB4 verify every comparison premise
of AD3: the lattices agree off $P,\bar P$, every other type functional
is the same, and all auxiliary maps and rational frames are unchanged.
AD3 supplies a unitary $\varepsilon$, equal to $1$ off that pair,
such that $y_ck\varepsilon$ realizes the raw neighbor. Set
\[
                         \eta=k\varepsilon k^{-1}. \tag{NB.18}
\]
Then $y_c\eta k=y_ck\varepsilon$. At every $F$-prime other than
$v_0=P\cap\mathcal O_F$, all its $E$-components are the identity,
so its projection to $E\otimes_FF_v$ is exactly $1$.
Conjugating by the same base point $k$ preserves this support and
unitarity. NB4 ensures that the new preferred class has the same
chart and therefore the very same base point $k$.

For the actual $g_0=(\gamma,h_0)$ of (NB.13), put
\[
 \kappa_0=y_{c'}^{-1}\operatorname{rat}(\gamma)^{-1}y_c\eta.
                                                               \tag{NB.19}
\]
The point realized by $y_{c'}\kappa_0k$ is
$\gamma^{-1}p_{\rm raw}=h_0r(c')$. Its coupling label equals that
of $r(c')$: NB4 proves equality of the raw and preferred labels,
and also invariance under the rational unitary $\gamma^{-1}$.
The exact level-reading equivalence in \texttt{ReadsK} therefore
places $\kappa_0$ in the chart level $U_\beta$; its character-reading
identity gives
\[
 \chi_\beta(\kappa_0)=\sigma(\chi_{\rm col}(g_0)).  \tag{NB.20}
\]
Here $\sigma$ is the supplied embedding of the coefficient field,
with the same identification used in CV and FC11. Set
$\kappa=\kappa_0^{-1}$. Rearranging (NB.19) and using the character
homomorphism proves
\[
 \operatorname{rat}(\gamma)y_{c'}=y_c\eta\kappa,
 \qquad \chi_\beta(\kappa)^{-1}
                   =\sigma(\chi_{\rm col}(g_0)).     \tag{NB.21}
\]
This one direct-splice proof supplies both the up factorization and
the conclusion of \texttt{move\_down\_adelic}, with the latter's
exact moving matrix. The source's down proof uses a reverse splice;
under the explicit realization input the direct splice above gives
the same required existential conclusion without changing its scalar.

\paragraph{NB7: the literal up naturality statement is now proved.}
Take the exact true-column data listed in the conventions, $i\le m$,
a chart class $c$, its selected prime $v$, and
$z\in\operatorname{LdG}(\operatorname{lat}c)$ with
$\mathrm{fv}_c(\operatorname{pr}_c z)=1$.
The coordinate identities of SF and FC identify
$\operatorname{lat}c=L_c$, $\operatorname{LdG}(\operatorname{lat}c)=H_c$,
and $\operatorname{PvG}(v)=P$. The identity \texttt{IdV} identifies
the base-field prime $(e_Vv)$ with $P\cap\mathcal O_F$, with its prescribed
orientation. Apply NB2 and NB4 to the quotient representative of
$\operatorname{pr}_c z$. They produce the source's actual up class
$c'=\mathrm{up}(v,c,\operatorname{pr}_c z)$ in the same chart.
Take its actual $\gamma$ from (NB.13), and $\eta,\kappa$ from NB6.
The support and adelic clauses are (NB.18)--(NB.21); the two lattice
clauses are (NB.14) with $q=T_u$; the functional clause is (NB.15).

Finally, for arbitrary $h$ and frames $s,s'$ related by
$e_Y(s)=\operatorname{atInf}(\gamma)e_Y(s')$, the literal
\texttt{IdY} and embedding identities give $\gamma^{-1}s=s'$
componentwise. The source's transport formula is consequently
\[
 \mathrm{upT}(v,c,\operatorname{pr}_c z)h(s)
       =\sigma(\chi_{\rm col}(g_0))h(s')
       =\chi_\beta(\kappa)^{-1}h(s').                \tag{NB.22}
\]
This checks every clause of \texttt{UpNatTieG}, and proves the actual
\texttt{UpChart.upNatTie}. The identity on frames uses
$\iota_\tau=\sigma\circ j_\tau$ at the supplied embeddings; no
complex conjugation, determinant factor or scalar normalization is
inserted into (NB.22).

\paragraph{NB8: the literal down naturality statement is now proved.}
For the same exact data, take any nonzero
$g\in\operatorname{Hom}_{k_P}(D_c/PD_c,k_P)$, which is the source's
\texttt{lk} dual under \texttt{IdL}. NB3 constructs the actual
$\mathrm{dnPt}$ using this $g$ and (NB.10), and NB4 puts its class
$c'=\mathrm{dn}(v,c,g)$ in the old chart. Use the actual $\gamma$
from (NB.13) and the supported $\eta,\kappa$ from NB6.
Equation (NB.21) supplies the exact adelic relation, (NB.16) supplies
the transported dual-lattice equality, and (NB.17) supplies the
functional comparison on every $x\in H_{c'}$ with
$\gamma x\in L_c+PH_c$. The same frame calculation as in NB7 yields
\[
                  \mathrm{dnT}(v,c,g)h(s)
                      =\chi_\beta(\kappa)^{-1}h(s'). \tag{NB.23}
\]
These are all the clauses of \texttt{DnNatTieG}, hence the actual
\texttt{DnMove.dnNatTieG\_cols}.

The conclusion closes the neighbor and coupling assertions formerly
retained in the Forms0 construction. Its scope remains conditional on
the literal supplied true-column laws and on initial adelic realization.
The latter is not supplied by a name for a record, by an equal-rank
argument, or by the local modifications just proved. NB makes no
additional claim about terminal-row incidence or the upstream
construction of \texttt{TC.Ok}, \texttt{Pins}, \texttt{Laws}, or
\texttt{Later}.
